\documentclass[10pt,a4paper]{amsart}
\usepackage[margin=19mm]{geometry}
\usepackage{amsmath,amssymb,mathtools}
\usepackage[scr=rsfs,cal=euler]{mathalfa}
\usepackage{xfrac}
\usepackage[unicode,hidelinks,bookmarks=false]{hyperref}
\newcommand{\ie}{i.e.\ }
\newcommand{\cf}{cf.\ }
\newcommand{\resp}{respectively\ }
\newcommand{\Subsection}{\subsection}
\providecommand{\nobreakdash}{\nobreak\hbox{-}}
\setlength{\emergencystretch}{2em}
\multlinegap0pt
\usepackage{amssymb}
\usepackage{mathtools}
\usepackage{xcolor}
\usepackage[shortlabels]{enumitem}
\newtheorem{lemma}{Lemma}
\def \R {\mathbb{R}}
\title{First integrals of dense hard-ball gases}
\author{Gleb Smirnov}
\date{Source: arXiv:2608.29694v1 (CC BY 4.0)}
\begin{document}
\maketitle

\noindent\emph{Source and licence.} First integrals of dense hard-ball gases. Version arXiv:2608.29694v1 (CC BY 4.0), distributed by arXiv under the Creative Commons Attribution 4.0 International licence, http://creativecommons.org/licenses/by/4.0/, as recorded in the arXiv licence metadata for this exact version and shown on the abstract page https://arxiv.org/abs/2608.29694v1. Full licence notice: \texttt{licenses/arxiv-2608.29694-CC-BY-4.0.txt}.\par\smallskip

\noindent\textit{Editorial scope.} Counted unit: the article's lemma lem:exchange (source0.tex lines 719--726). The article's complete original proof of it is delivered with it (source0.tex lines 727--839, one \texttt{proof} environment of 113 lines). No mathematics is written, repaired or re-derived: the statement and the proof are the article's own lines, in the article's own notation, and no retained line is substituted or re-flowed.  No further source-local context is needed: every cross-reference of every retained line resolves inside the two delivered spans.  Nothing was omitted from any retained span, so the package carries no omission record.  No retained line contains a citation, so the wrapper carries no bibliography and no bibliography entry.  No retained line contains a figure environment, an image-inclusion command or drawing code, so no image is delivered and the package declares no asset dependency.  Editorial formatting only, in addition: an AMS \texttt{amsart} preamble that restates the article's own declarations and macros used by the retained lines, namely the environments \texttt{lemma} on separate counters, together with the standard packages those lines need.  The retained environments print in this document as Lemma 1 in order of appearance, whereas the article prints the same items with the numbering of its own full document.  The complete native source of the article as distributed in the arXiv e-print for this exact version is bundled unchanged as \texttt{sources/208879/source0.tex}.
\begin{lemma}[Pair symmetry]\label{lem:exchange}
Suppose that \(ij\) is an edge of \(\Gamma_\gamma\). Let
\(\tau_{ij}\) exchange the \(i\)-th and \(j\)-th balls. Then:
\[
f(\tau_{ij}q,\tau_{ij}p)=f(q,p),
\qquad q,p\in\R^{dN}.
\]
\end{lemma}

\begin{proof}
Set:
\[
\mathcal C_{ij}
=
\left\{
q\in\R^{dN}:\|q_i-q_j\|=2\rho
\right\}.
\]
Since \(ij\) is an edge of \(\Gamma_\gamma\),
\[
\mathcal C_{ij}\cap\overline{\Omega}_\gamma
\]
contains a nonempty relatively open subset of \(\mathcal C_{ij}\). For \(q\in\mathcal C_{ij}\), set:
\[
n=\frac{q_i-q_j}{2\rho}.
\]
On this subset the collision law gives:
\[
\begin{aligned}
p_i'&=p_i-\langle p_i-p_j,n\rangle n,\\
p_j'&=p_j+\langle p_i-p_j,n\rangle n,
\end{aligned}
\]
with all other momenta unchanged; hence:
\[
f(q,p')=f(q,p)
\]
on a nonempty relatively open subset of
\(\mathcal C_{ij}\times\R^{dN}\). Since both sides are polynomial and
\(\mathcal C_{ij}\) is connected, it 
follows that:
\[
f(q,p')=f(q,p)\quad \text{on the whole $\mathcal C_{ij}$.}
\]
Introduce the coordinates:
\[
x=\frac{q_i-q_j}{\sqrt2},
\qquad
v=\frac{p_i-p_j}{\sqrt2},
\]
and complete them to orthogonal 
coordinates:
\[
(x,y)\in\R^d\times\R^{dN-d},
\qquad
(v,w)\in\R^d\times\R^{dN-d}.
\]
Then the collision surface is \(\|x\|=a\), \(a = \sqrt{2} \rho\), and the collision law becomes:
\[
v\longmapsto
r_xv
=
v-2\frac{\langle v,x\rangle}{a^2}x.
\]
Thus:
\begin{equation}\label{eq:ball-reflection}
f(x,y,v,w)=f(x,y,r_xv,w),
\qquad \|x\|=a.
\end{equation}
The global free-flight equation 
allows us to consider the 
following formal billiard. The variable \(x\) moves freely inside the 
ball \(\|x\| \le a\), 
and is reflected at the sphere 
\(\|x\|=a\). The remaining position 
variables move freely:
\[
y(t)=y+tw,
\]
while \(w\) remains fixed. This need not be an actual trajectory of the
hard-ball system: other balls may overlap or leave the box. Nevertheless, \(f\) is constant along it. 
\smallskip%

Consider an initial state \((x,v)\) 
with \(\|x\|< a, v\neq 0\). Suppose that \(\|x\|<a\) and \(x,v\) are not 
collinear; then the \(x\)-motion stays in the two-dimensional plane spanned by \(x\) and
\(v\), and is an ordinary 
disc billiard. For a dense set of states \((x,v)\) the billiard trajectory 
is periodic. For some 
\(T = T(x,v) > 0\),
the trajectory returns to \((x,v)\) 
after time \(T\). Consequently,
\[
 f(x,y,v,w)
=
f(x,y+Tw,v,w).
\]
For fixed \(x,v,w\), the function \(y\mapsto f(x,y,v,w)\) is polynomial. If \(w\ne0\), the polynomial \(y\to f(x,y,v,w)\) has a nonzero
period; hence:
\begin{equation}\label{eq:period}
f(x,y+sw,v,w)=f(x,y,v,w) \quad s\in\R.
\end{equation}
Since periodic states are dense and \(f\) is polynomial, this identity
holds for all \(x,v,y,w\).
\smallskip%

Take a state \((x,v)\) such that the 
trajectory hits \(-x, -v\) after 
some time \(T\). Then:
\[
f(x,y,v,w) = f(-x,y + T w,-v,w).
\]
From \eqref{eq:period}: 
\[
f(x,y,v,w) = f(-x,y,-v,w).
\]
Such states are dense, so the identity holds everywhere. Finally,
\[
(x,y,v,w)\longmapsto(-x,y,-v,w)
\]
is precisely the simultaneous exchange of balls \(i\) and \(j\). This completes the proof. \qed
\end{proof}

\end{document}
